\section{Collisionless closures}
\label{sec:collisionless}

In hot, dilute plasmas the mean free path exceeds the system size. Nothing
relaxes $f$ toward a Maxwellian, and Chapman--Enskog has no small
parameter. What takes the place of collisions is \emph{phase mixing}:
particles streaming at different speeds smear out perturbations, and
particles resonant with a wave ($v=\omega/k$) exchange energy with it
(Landau damping \cite{Landau46}). This section derives the closure that
builds phase mixing into a fluid model, following Hammett and Perkins
\cite{HP90,HDP92}. The derivation is the one in the companion note, with
the same conventions; the closure at general order
(Section~\ref{sec:general-order}) and the error figures are new here.

\subsection{Setup and conventions}

One species, one spatial dimension $z$ (along the magnetic field), no
collisions, linearized about
\[
  f_0(v)=\frac{n_0}{\sqrt{2\pi}\,\vt}\exp\Bigl(-\frac{v^2}{2\vt^2}\Bigr),
  \qquad \vt=\sqrt{T_0/m}.
\]
With acceleration $a=(e/m)E$ the perturbation obeys
\begin{equation}
  \partial_tf_1+v\,\partial_zf_1+a\,\partial_vf_0=0 .
  \label{eq:vlasov}
\end{equation}
Derivations use units $n_0=m=\vt=1$, in which this is \eqref{eq:linkin}
with $\Lambda=0$; dimensional forms are restored for results. Fourier modes
are $\propto e^{ikz-i\omega t}$, and
\begin{equation}
  \zeta=\frac{\omega}{\sqrt2\,|k|\,\vt}.
  \label{eq:zeta}
\end{equation}
The Hilbert transform is taken with the convention
\begin{equation}
  \Hil[g](z)=\frac1\pi\,\PV\int_{-\infty}^{\infty}\frac{g(z')}{z-z'}\,\dd z',
  \qquad
  \widehat{\Hil[g]}(k)=-i\,\sgn(k)\,\hat g(k).
  \label{eq:hilbert}
\end{equation}

\subsection{The kinetic response and the origin of $\sgn(k)$}

In Fourier space \eqref{eq:vlasov} gives, using $\partial_vf_0=-vf_0$,
$f_{1,k}=-i\,a\,v\,f_0/(kv-\omega)$, so the density perturbation is
\[
  n_k=-\frac{ia}{k}\int\frac{v\,f_0}{v-\omega/k}\,\dd v
     =-\frac{ia}{k}\Bigl[1+\frac\omega k\int\frac{f_0}{v-\omega/k}\,\dd v\Bigr].
\]
The integral is defined for $\operatorname{Im}\omega>0$ (causality) and
continued analytically. It is the plasma dispersion function
\begin{equation}
  Z(\zeta)=\frac1{\sqrt\pi}\int_{-\infty}^{\infty}\frac{e^{-x^2}}{x-\zeta}\,\dd x,
  \qquad\operatorname{Im}\zeta>0,
  \label{eq:Z}
\end{equation}
the Hilbert transform, in velocity space, of a Gaussian. For real $\zeta$
the Plemelj formula
$1/(x-\zeta-i0)=\PV\,1/(x-\zeta)+i\pi\,\delta(x-\zeta)$ separates a
principal-value part from the contribution of the resonant particles.

The pole is at $v=\omega/k$. For $\operatorname{Im}\omega>0$ it lies above
the real $v$ axis when $k>0$ and below when $k<0$. Substituting $x\to-x$ in
the second case shows that both signs give the same function of $\zeta$ as
defined in \eqref{eq:zeta}, with $|k|$ and not $k$:
\begin{equation}
  n_k=\frac{a}{ik}\,R(\zeta),\qquad R(\zeta)=1+\zeta Z(\zeta).
  \label{eq:R}
\end{equation}
This is the source of the non-analytic dependence on $k$. The expansions
needed are
\begin{align}
  R(\zeta)&=1+i\sqrt\pi\,\zeta-2\zeta^2+O(\zeta^3), &\zeta&\to0,
     \label{eq:Rsmall}\\
  R(\zeta)&=-\frac1{2\zeta^2}-\frac3{4\zeta^4}-\frac{15}{8\zeta^6}-\dots,
     &|\zeta|&\to\infty. \label{eq:Rlarge}
\end{align}
The term $i\sqrt\pi\,\zeta$ is the resonant contribution,
$Z(0)=i\sqrt\pi$. On the real axis $|R|\le1$.

\subsection{Linearized moment equations}

Define $n=\int f_1\,\dd v$, $u=\int vf_1\,\dd v$, $p=\int v^2f_1\,\dd v$,
$T=p-n$, the heat flux $q=\int v^3f_1\,\dd v-3u$, and the deviation of the
fourth moment from its Maxwellian value,
$\delta r=\int v^4f_1\,\dd v-6p+3n$. In the Hermite variables of
Section~\ref{sec:hermite}, $(n,u,T,q,\delta r)
=(a_0,a_1,\sqrt2a_2,\sqrt6a_3,\sqrt{24}a_4)$. Moments of \eqref{eq:vlasov}
give
\begin{align}
  \partial_tn+\partial_zu&=0, \label{eq:n}\\
  \partial_tu+\partial_z(n+T)&=a, \label{eq:u}\\
  \partial_tT+2\,\partial_zu+\partial_zq&=0, \label{eq:T}\\
  \partial_tq+3\,\partial_zT+\partial_z\delta r&=0. \label{eq:q}
\end{align}
Each is exact. In Fourier space write
$\omega/k=\sqrt2\,\sgn(k)\,\zeta$.

\subsection{Three-moment closure}

Keep \eqref{eq:n}--\eqref{eq:T} and close $q$. A diffusive closure is
$q_k=-\chi\,ik\,T_k$. With collisions $\chi\sim\vt^2/\nu$. Without them the
only available rate is the phase-mixing rate $|k|\vt$, so
$\chi=\chi_1\sqrt2\,\vt/|k|$ with $\chi_1$ a pure number:
\begin{equation}
  q_k=-i\,\sgn(k)\,\sqrt2\,\chi_1\,T_k .
  \label{eq:closure3}
\end{equation}
Equation \eqref{eq:T} then gives $T_k=2\zeta\,n_k/(\zeta+i\chi_1)$, and
\eqref{eq:u} gives $n_k=(a/ik)\,R_3$ with
\begin{equation}
  R_3(\zeta)=\frac{\chi_1-i\zeta}
     {\chi_1-3i\zeta-2\chi_1\zeta^2+2i\zeta^3}.
  \label{eq:R3}
\end{equation}
For large $\zeta$, $R_3=-1/(2\zeta^2)-3/(4\zeta^4)+O(\zeta^{-5})$, which
agrees with \eqref{eq:Rlarge} for any $\chi_1$ because the retained moment
equations are exact. For small $\zeta$,
$R_3=1+2i\zeta/\chi_1+O(\zeta^2)$, and comparison with \eqref{eq:Rsmall}
fixes
\begin{equation}
  \chi_1=\frac2{\sqrt\pi}.
\end{equation}
Restoring units,
\begin{equation}
  q_k=-n_0\,\chi_1\frac{\sqrt2\,\vt}{|k|}\,ik\,T_k
     =-i\,\sgn(k)\;n_0\sqrt{\frac8\pi}\;\vt\;T_k .
  \label{eq:hp3}
\end{equation}
By \eqref{eq:hilbert} the multiplier $-i\,\sgn(k)$ is the Hilbert
transform, so in real space
\begin{equation}
  q(z)=n_0\sqrt{\frac8\pi}\;\vt\;\Hil[T](z)
      =-n_0\,\chi_1\frac{\sqrt2\,\vt}{\pi}
        \int_0^\infty\frac{T(z+z')-T(z-z')}{z'}\,\dd z' .
  \label{eq:real3}
\end{equation}
The kernel carries no length scale, because a collisionless plasma has no
mean free path to supply one.

Compare the two heat-flux closures:
\begin{equation}
  \text{collisional:}\quad q_k=-\kappa\,ik\,T_k,
  \qquad\qquad
  \text{collisionless:}\quad
  q_k=-n_0\sqrt{\frac8\pi}\;\vt\;i\,\sgn(k)\,T_k .
\end{equation}
The first damps temperature perturbations at a rate $\propto k^2$, the
second at a rate $\propto|k|\vt$. In real space the first is a derivative
and the second a nonlocal operator. Both are dissipative: with
\eqref{eq:hp3},
$\int q\,\partial_zT\,\dd z=-n_0\sqrt{8/\pi}\,\vt\sum_k|k|\,|T_k|^2\le0$.

\subsection{Four-moment closure}

Now keep \eqref{eq:q} as well and close $\delta r$ in terms of the two
highest retained moments. Parity and dimensions allow a local term in $T$
and a phase-mixing term in $q$:
\begin{equation}
  \delta r_k=\beta_1\,T_k-i\,\sgn(k)\,\sqrt2\,D_1\,q_k .
  \label{eq:closure4}
\end{equation}
The $T$ term is not dissipative; it replaces the coefficient $3$ in
\eqref{eq:q} by $b\equiv3+\beta_1$. Eliminating upward through the
hierarchy,
\begin{align*}
  \text{from \eqref{eq:q}:}\quad
    q_k&=\frac{b\,T_k}{\sqrt2\,\sgn(k)\,(\zeta+iD_1)},\\[4pt]
  \text{from \eqref{eq:T}:}\quad
    T_k&=\frac{4\zeta(\zeta+iD_1)}{2\zeta(\zeta+iD_1)-b}\;n_k,\\[4pt]
  \text{from \eqref{eq:u}:}\quad
    n_k&=\frac{a}{ik}\,R_4(\zeta),
\end{align*}
with
\begin{equation}
  R_4(\zeta)=\frac{b-2iD_1\zeta-2\zeta^2}
     {b-6iD_1\zeta-(2b+6)\zeta^2+4iD_1\zeta^3+4\zeta^4}.
  \label{eq:R4}
\end{equation}
Every factor of $\sgn(k)$ cancels except through $\zeta$.

\paragraph{Matching at small $\zeta$.} With $\alpha=2iD_1/b$,
\[
  R_4(\zeta)=1+2\alpha\,\zeta
    +\Bigl(2+\frac4b+6\alpha^2\Bigr)\zeta^2+O(\zeta^3).
\]
Comparing with \eqref{eq:Rsmall}:
\begin{align*}
  O(\zeta):&\quad2\alpha=i\sqrt\pi
    \;\Longrightarrow\;D_1=\frac{\sqrt\pi}{4}\,b,\\
  O(\zeta^2):&\quad2+\frac4b-\frac{3\pi}{2}=-2
    \;\Longrightarrow\;b=\frac{8}{3\pi-8}.
\end{align*}
Therefore
\begin{equation}
  D_1=\frac{2\sqrt\pi}{3\pi-8}\approx2.49,
  \qquad
  \beta_1=\frac{32-9\pi}{3\pi-8}\approx2.62 .
  \label{eq:coeffs4}
\end{equation}

\paragraph{Behaviour at large $\zeta$.} From \eqref{eq:R4},
\[
  \frac1{R_4}=-2\zeta^2+3+\frac{b}{\zeta^2}+O(\zeta^{-3}),
  \qquad\text{whereas}\qquad
  \frac1R=-2\zeta^2+3+\frac3{\zeta^2}+\dots
\]
So $R_4$ reproduces the first two terms of \eqref{eq:Rlarge} for any $D_1$
and $\beta_1$, but the $\zeta^{-6}$ term is correct only if $\beta_1=0$.
The closure spends its two parameters on the low-frequency side and gives
up one order of high-frequency accuracy to do so.

\paragraph{Real space.} Restoring units,
\begin{equation}
  \delta r(z)=\beta_1\,n_0\vt^2\,T(z)+\sqrt2\,D_1\,\vt\,\Hil[q](z),
\end{equation}
to be used in
$\partial_tq+3n_0\vt^2\,\partial_zT+\partial_z\delta r=0$. The Hilbert
transform now acts on the heat flux, which has its own evolution equation.

\subsection{How accurate the closures are}

\begin{center}
\renewcommand{\arraystretch}{1.5}
\begin{tabular}{lll}
\toprule
  & Three-moment & Four-moment \\
\midrule
  Closed quantity & $q$ & $\delta r$ \\
  Closure & $q=n_0\sqrt2\,\chi_1\vt\,\Hil[T]$
          & $\delta r=\beta_1n_0\vt^2T+\sqrt2\,D_1\vt\,\Hil[q]$ \\
  Coefficients & $\chi_1=2/\sqrt\pi$
          & $D_1=\dfrac{2\sqrt\pi}{3\pi-8}$,\;
            $\beta_1=\dfrac{32-9\pi}{3\pi-8}$ \\
  Matched as $\zeta\to0$ & through $O(\zeta)$ & through $O(\zeta^2)$ \\
  Matched as $\zeta\to\infty$ & through $O(\zeta^{-4})$
          & through $O(\zeta^{-4})$ \\
  $\sup_{\zeta\in\mathbb R}|R_n-R|$ & $0.34$ at $\zeta\approx1.0$
          & $0.059$ at $\zeta\approx1.3$ \\
\bottomrule
\end{tabular}
\end{center}
The last row was computed for this note; Figure~\ref{fig:response} shows
the functions. Both closures are exact at the two ends of the frequency
axis and are rational (Pad\'e-type) approximants in between. The
three-moment closure is poor at phase velocities near the thermal speed,
where most particles are resonant; the four-moment closure is within six
percent of the peak response everywhere on the real axis.

\begin{figure}[ht]
\centering
\begin{tikzpicture}
\begin{axis}[width=0.48\textwidth,height=5.6cm,xlabel={$\zeta$},
  title={$\operatorname{Re}R$},xmin=0,xmax=4,grid=major,
  legend style={font=\footnotesize,at={(0.97,0.97)},anchor=north east}]
\addplot[black,thick] table[x=zeta,y=ReR]{\rundir/response.dat};
\addplot[blue,dashed,thick] table[x=zeta,y=ReR3]{\rundir/response.dat};
\addplot[red,densely dotted,very thick] table[x=zeta,y=ReR4]{\rundir/response.dat};
\legend{exact,3-moment,4-moment}
\end{axis}
\end{tikzpicture}\hfill
\begin{tikzpicture}
\begin{axis}[width=0.48\textwidth,height=5.6cm,xlabel={$\zeta$},
  title={$\operatorname{Im}R$},xmin=0,xmax=4,grid=major]
\addplot[black,thick] table[x=zeta,y=ImR]{\rundir/response.dat};
\addplot[blue,dashed,thick] table[x=zeta,y=ImR3]{\rundir/response.dat};
\addplot[red,densely dotted,very thick] table[x=zeta,y=ImR4]{\rundir/response.dat};
\end{axis}
\end{tikzpicture}
\caption{The collisionless density response $R(\zeta)=1+\zeta Z(\zeta)$ on
the real axis and its three- and four-moment approximants \eqref{eq:R3},
\eqref{eq:R4}. Data: \texttt{\rundir/response.dat}.}
\label{fig:response}
\end{figure}

\subsection{The closure at general order}
\label{sec:general-order}

The three-moment closure has a simple description that extends to any
order. In the Hermite hierarchy \eqref{eq:hermitehier} with $\Lambda=0$,
the forced response at frequency $\omega$ has coefficients, for $n\ge1$,
\[
  \hat a_n\;\propto\;\frac{W_n(z)}{\sqrt{n!}},\qquad
  W_n(z)=\Bigl\langle\frac{\He_n(v)}{v-z}\Bigr\rangle,\qquad
  z=\frac{\omega}{k},
\]
and the Hermite recurrence gives
$W_{n+1}=z\,W_n-n\,W_{n-1}+\delta_{n0}$, with
$W_0(0)=i\sqrt{\pi/2}\,\sgn(k)$ and $W_1(0)=1$. At zero frequency the
ratio of consecutive coefficients is therefore an explicit number:
\begin{equation}
  \hat a_{N+1}=-i\,\sgn(k)\,c_N\,\hat a_N,
  \qquad
  c_N=\frac{N!!}{(N-1)!!\,\sqrt{N+1}}\times
  \begin{cases}\sqrt{2/\pi},&N\text{ even},\\[2pt]
               \sqrt{\pi/2},&N\text{ odd}.\end{cases}
  \label{eq:cN}
\end{equation}
Equivalently $c_Nc_{N+1}=\sqrt{(N+1)/(N+2)}$ with
$c_2=4/\sqrt{6\pi}$. In terms of the continued fraction \eqref{eq:cf},
$c_N=\sqrt{N+1}\,|k|/\tau_{N+1}(0)$: closing with \eqref{eq:cN} replaces
the tail of the continued fraction by its value at zero frequency.

\begin{center}
\begin{tabular}{lrrrrrrr}
\toprule
$N$ & 2 & 3 & 4 & 5 & 6 & 7 & 8 \\
\midrule
$c_N$ & 0.9213 & 0.9400 & 0.9515 & 0.9594 & 0.9650 & 0.9693 & 0.9727 \\
\bottomrule
\end{tabular}
\end{center}
(The formula was checked against direct quadrature of $W_n$ at
$z=10^{-4}i$; agreement is to $2\times10^{-4}$, the accuracy of the
quadrature.)

Three observations.
\begin{itemize}
  \item $N=2$ is Hammett--Perkins:
    $q=\sqrt6\,a_3$ and $T=\sqrt2\,a_2$ give
    $q_k=-i\,\sgn(k)\sqrt3\,c_2\,T_k=-i\,\sgn(k)\sqrt{8/\pi}\,T_k$.
  \item $N=3$ gives $\delta r_k=-i\,\sgn(k)\,\tfrac32\sqrt{\pi/2}\;q_k$,
    which is \eqref{eq:closure4} with $\beta_1=0$ and
    $D_1=\tfrac{3}{4}\sqrt\pi$. It matches $R$ through $O(\zeta)$ only. The
    Hammett--Perkins four-moment closure uses the additional term
    $\beta_1T$ to match $O(\zeta^2)$ as well, at the cost described above.
  \item $c_N\to1$ as $N\to\infty$. The limiting closure
    $\hat a_{N+1}=-i\,\sgn(k)\,\hat a_N$ is an outgoing-wave condition:
    free streaming carries perturbations toward high $n$, and the closure
    lets them leave instead of reflecting at the truncation. This is the
    absorber that Section~\ref{sec:moment} said a truncated system lacks.
\end{itemize}

\subsection{Caveats}
\label{sec:hp-caveats}

\begin{itemize}
  \item \textbf{Thermal-speed convention.} This note and Hammett--Perkins
    use $\vt=\sqrt{T_0/m}$. Hunana et al.\ \cite{Hunana19} use
    $v_{\mathrm{th}}=\sqrt{2T_0/m}$, in which the three-moment closure
    reads $q_k=-i\,(2/\sqrt\pi)\,n_0v_{\mathrm{th}}\,\sgn(k)\,T_k$ and the
    local coefficient of the four-moment closure appears a factor of two
    smaller.
  \item \textbf{Heat-flux convention.} Here $q$ is the third central
    moment. The heat-flux manuscript uses the energy flux, which is half of
    it in one dimension, so $\sqrt{8/\pi}$ becomes $\sqrt{2/\pi}$ there
    (Appendix~\ref{app:conv}).
  \item \textbf{Sign of the Hilbert transform.} With \eqref{eq:hilbert} and
    modes $e^{+ikz}$, $i\,\sgn(k)\leftrightarrow-\Hil$. Several
    Landau-fluid papers define the opposite sign. The integral form
    \eqref{eq:real3} is convention-free.
  \item \textbf{Known inaccuracies.} According to \cite{Hunana19}, the
    three-moment closure gives $q\sim\zeta^{-2}$ at large $\zeta$ where the
    kinetic result is $\zeta^{-3}$, and so tends to overestimate Landau
    damping. The closures are local in time.
  \item \textbf{Linearity.} The matching is linear, about a Maxwellian.
    Collisionless phase mixing is reversible: the perturbation is still in
    the distribution function, at fine velocity scales, and a second
    perturbation can bring it back as a plasma echo \cite{Gould67}. A
    closure that damps cannot do that. Nor does it describe particle
    trapping, which stops Landau damping after a bounce time at finite
    amplitude.
  \item \textbf{Domain.} $\Hil$ is defined on the whole line or a periodic
    domain. Bounded or sheath-terminated field lines need a separate choice
    of extension.
\end{itemize}

\subsection{Extensions used in practice}

\begin{itemize}
  \item \textbf{Three dimensions and magnetization.} The parallel and
    perpendicular pressures are evolved separately, with one closure for
    each parallel heat flux \cite{Snyder97}. Gyrofluid models add closures
    for finite-gyroradius averaging and for drift resonances
    \cite{DorlandHammett93}.
  \item \textbf{Collisions.} The collisionless and collisional rates are
    usually combined by adding them, which is the harmonic interpolant of
    Section~\ref{sec:between}. The exact answer for model collision
    operators is known there and can be compared.
  \item \textbf{Real-space evaluation.} Along curved or bounded field lines
    a Fourier transform is unavailable. Since
    \begin{equation}
      \frac1{|k|}=\frac1\pi\int_0^\infty\frac{m^{-1/2}}{m+k^2}\,\dd m,
      \label{eq:lorentz}
    \end{equation}
    the operator $1/|k|$ can be approximated over a band of wavenumbers by
    a short sum of terms $w_j/(m_j+k^2)$, each of which is a tridiagonal
    solve in real space \cite{DJU14}. Equation \eqref{eq:lorentz} says that
    the collisionless closure is the limiting member of the same family of
    screened-diffusion closures as the collisional nonlocal models of
    Section~\ref{sec:plasma}.
  \item \textbf{Relaxation closures in ten-moment models} replace
    $|k|$ by a fixed wavenumber $k_0$, giving a local term that relaxes the
    pressure tensor toward isotropy at a rate $\vt|k_0|$ \cite{Wang15}.
\end{itemize}
